\chapter{Options, forwards and absence of arbitrage}
\label{cap:opzioni}

In this chapter we build the basic language: what an option is, what the discount factor and the forward price are,
and why the principle of \emph{absence of arbitrage} is enough on its own to derive exact relations between prices
(put--call parity) without any model. At the end we will see why the price of a derivative is written as an
\emph{expected value}, and under which probability.

\section{What an option is}

\begin{definizione}[European call and put]
A \textbf{European call} with strike $K$ and expiry $T$ gives the right, but not the obligation, to buy the
underlying at the price $K$ on date $T$. A \textbf{European put} gives the right to sell it at $K$ on $T$.
\end{definizione}

The holder of the call exercises it only if it pays to do so, that is, if $S_T>K$: she buys at $K$ something worth
$S_T$ and gains $S_T-K$. Otherwise she lets the option expire and receives zero. The \textbf{payoff} (the cash flow at
expiry) is therefore
\begin{equation}
\text{call: } (S_T-K)^+ ,\qquad \text{put: } (K-S_T)^+ ,\qquad\text{where } y^+:=\max(y,0).
\end{equation}
In compact form, with $\theta=+1$ for the call and $\theta=-1$ for the put: $(\theta(S_T-K))^+$.

\begin{center}
\begin{tikzpicture}[scale=0.9,>=Stealth]
\draw[->] (0,0)--(5,0) node[right]{$S_T$}; \draw[->] (0,0)--(0,3) node[above]{payoff};
\draw[very thick,bluidea] (0,0)--(2.2,0)--(4.8,2.6); \node[below] at (2.2,0){$K$};
\node[bluidea] at (3.2,2.3){call};
\begin{scope}[xshift=7cm]
\draw[->] (0,0)--(5,0) node[right]{$S_T$}; \draw[->] (0,0)--(0,3) node[above]{payoff};
\draw[very thick,rossoatt] (0,2.2)--(2.2,0)--(4.8,0); \node[below] at (2.2,0){$K$};
\node[rossoatt] at (1.6,2.0){put};
\end{scope}
\end{tikzpicture}
\end{center}

The SPX index options listed on Cboe are \textbf{European} (exercisable only at expiry) and \textbf{cash-settled}
(the index is not delivered: the difference is paid). This is a huge advantage for us: all the ``European''
relations we will see hold exactly. Options on the SPY ETF, by contrast, are \emph{American} (exercisable at any
time), and for them many equalities become mere inequalities.

\paragraph{Basic terminology.} A call is \emph{in the money} (ITM) if $S>K$, \emph{at the money} (ATM) if
$S\approx K$, \emph{out of the money} (OTM) if $S<K$; for the put it is the other way round. The \emph{intrinsic
value} is the payoff if the option were exercised immediately, $(\theta(S-K))^+$; the \emph{time value} is the price
minus the intrinsic value. An OTM option has zero intrinsic value: its price is all time value. This observation will
come back several times (Chapters~\ref{cap:dati} and~\ref{cap:black}).

\section{The time value of money: the discount factor}

A dollar today is worth more than a dollar a year from now, because it can be invested. The \textbf{discount factor}
\begin{equation}
\DF = P(0,T) = \text{today's price of a security that pays 1 at } T
\end{equation}
(a \emph{zero-coupon bond}) summarizes this effect. The continuously compounded \emph{zero} rate is defined by
$\DF=e^{-rT}$, that is
\begin{equation}\label{eq:tasso-da-DF}
r = -\frac{\ln \DF}{T}.
\end{equation}
With positive rates, $\DF\in(0,1]$. In the project the discount factor is not taken from an external curve: it is
\emph{inferred from the option prices} (Chapter~\ref{cap:dati}).

\section{The forward price}

A \textbf{forward contract} with delivery price $K$ obliges its holder to buy the underlying at $T$ paying $K$. Its
payoff is $S_T-K$ (it can be negative: it is an obligation). The \textbf{forward price} $F=F(0,T)$ is the value of
$K$ that makes today's value of the contract zero.

\paragraph{How it is derived (cash and carry).} Suppose the underlying pays no dividends. I can replicate the payoff
$S_T-K$ as follows: today I buy the underlying at $S_0$ and borrow $K\cdot\DF$ (that is, I sell $K$ zero-coupon
bonds). At $T$ I hold the underlying, worth $S_T$, and must repay $K$: in total $S_T-K$. Two strategies with the same
payoff must have the same price today (otherwise I buy the cheaper one, sell the dearer one and pocket a riskless
profit). Hence
\begin{equation}
\text{forward value} = S_0-K\,\DF \quad\Longrightarrow\quad F=\frac{S_0}{\DF}=S_0e^{rT}.
\end{equation}
If the underlying pays dividends (or earns something when lent out, the \emph{repo}), whoever holds the underlying
collects them and whoever holds the forward does not; the forward is lower:
\begin{equation}\label{eq:forward-dividendi}
F=\frac{S_0-\mathrm{PV}(\text{dividends})}{\DF},
\end{equation}
or, with a continuous yield $q$, $F=S_0e^{(r-q)T}$.

\begin{idea}
The forward $F$ already contains \emph{all} of the ``carry'': rates, dividends, stock-borrow cost. Today's value of a
forward with delivery price $K$ is
\begin{equation}\label{eq:valore-forward}
\text{value} = \DF\,(F-K).
\end{equation}
Indeed, at expiry the contract pays $S_T-K$; the forward with price $F$ is worth zero, so the contract with price $K$
is worth as much as receiving the certain difference $F-K$ at $T$, that is, $\DF\,(F-K)$.
\end{idea}

\section{Absence of arbitrage and put--call parity}

\begin{definizione}[Arbitrage]
An \textbf{arbitrage} is a strategy that costs zero (or less) today, never loses, and makes money in some scenario.
In a liquid market arbitrages are eliminated instantly; one therefore assumes that they \emph{do not exist}.
\end{definizione}

The most important consequence is the \emph{law of one price}: two portfolios with the same payoff in every scenario
have the same price today.

\begin{proposizione}[Put--call parity]\label{prop:parita}
For European calls and puts with the same strike $K$ and the same expiry $T$:
\begin{equation}\label{eq:parita}
C-P=\DF\,(F-K).
\end{equation}
\end{proposizione}

\begin{proof}[Step-by-step reasoning]\leavevmode
\begin{enumerate}
\item At expiry, for any value of $S_T$: $(S_T-K)^+-(K-S_T)^+=S_T-K$. (If $S_T>K$ the first term equals $S_T-K$ and
the second is zero; if $S_T<K$ the first is zero and the second is $K-S_T$, which with the minus sign again gives
$S_T-K$.)
\item Therefore the portfolio ``long one call, short one put'' has \emph{the same payoff} as a forward with delivery
price $K$.
\item By the law of one price the values today coincide; with \eqref{eq:valore-forward},
$C-P=\DF\,(F-K)$.\qedhere
\end{enumerate}
\end{proof}

\begin{perche}[title={Why is it so important?}]
Because it \textbf{uses no model}: there is no need to know how the price moves or how much volatility there is, nor
to know rates and dividends. It holds for any European option. The project uses it in three ways:
\begin{enumerate}
\item in reverse, to \emph{measure} $F$ and $\DF$ from the market prices of calls and puts
(Chapter~\ref{cap:dati});
\item as a correctness test for the code (if the pricer violates parity, there is a bug);
\item to obtain a call from a put already computed (useful when computing the put is numerically more stable, as in
Heston, Chapter~\ref{cap:heston}).
\end{enumerate}
\end{perche}

\subsection{Price bounds and static arbitrage in \texorpdfstring{$K$}{K}}
Other ``model-free'' constraints follow from the same principle. Denote by $c=C/\DF$ and $p=P/\DF$ the
\emph{undiscounted} prices.
\begin{itemize}
\item \textbf{Bounds:} $(F-K)^+\le c\le F$ and $(K-F)^+\le p\le K$. The call is worth at least as much as the forward
(I would never exercise it at a loss) and no more than the underlying itself.
\item \textbf{Monotonicity:} the call is decreasing in $K$ (a higher strike pays less in every scenario), the put
increasing.
\item \textbf{Vertical spread:} for $K_1<K_2$, $0\le C(K_1)-C(K_2)\le\DF\,(K_2-K_1)$: the portfolio ``long the $K_1$
call, short the $K_2$ call'' pays between $0$ and $K_2-K_1$.
\item \textbf{Butterfly (convexity):} $C$ is convex in $K$. The combination ``long $C(K-h)$, long $C(K+h)$, short
$2C(K)$'' pays a non-negative ``tent'', so it must cost $\ge0$.
\end{itemize}

\begin{idea}[title={Key idea: from prices to the density (Breeden--Litzenberger)}]
Differentiating the call price twice with respect to the strike gives the probability density of the underlying at
expiry (under the pricing measure we will meet shortly):
\begin{equation}\label{eq:breeden}
\frac{\partial^2 C}{\partial K^2}=\DF\cdot f_{S_T}(K).
\end{equation}
The reason: $\partial_K(S_T-K)^+=-\ind\{S_T>K\}$ and $\partial_K^2(S_T-K)^+=\delta(S_T-K)$ (a Dirac delta). Convexity
in $K$ therefore simply means ``non-negative density''. A set of prices that violates it implies negative
probabilities, that is, an arbitrage (``butterfly arbitrage'').
\end{idea}

\section{Why the price is an expected value}

\subsection{A one-step example}
Consider a world with two scenarios: today the underlying is worth $S_0=100$; one year from now it is worth $120$
(``up'') or $80$ (``down''). For simplicity, rates are zero. A call with strike $100$ pays $20$ or $0$.

\emph{Replication.} I look for $\Delta$ units of the underlying and $B$ dollars in cash such that the portfolio pays
like the call: $120\Delta+B=20$ and $80\Delta+B=0$. Subtracting: $40\Delta=20$, so $\Delta=0.5$ and $B=-40$. The
portfolio costs $0.5\cdot100-40=10$ today. Then the call must be worth $10$: otherwise there is an arbitrage.

\emph{Crucial observation.} We never used the ``real'' probability that the price goes up! Now I define $q$ as the
probability that makes the underlying a \emph{martingale} (expected future value equal to today's value):
$120q+80(1-q)=100\Rightarrow q=\tfrac12$. Then
\[
\E^q[\text{payoff}]=\tfrac12\cdot20+\tfrac12\cdot0=10 .
\]
The same number. This is no coincidence: the no-arbitrage price is always the (discounted) expected value of the
payoff under the probability that makes discounted prices martingales. This probability is called the
\textbf{risk-neutral measure} (or martingale measure). It is \emph{not} a forecast: it is a computational tool that
encodes the fact that the derivative can be replicated.

\subsection{Numéraire and forward measure}
The argument generalizes (the fundamental theorem of asset pricing). One chooses an always-positive reference
security, the \textbf{numéraire} $N_t$; there exists a measure $\Q^N$ under which \emph{every price divided by the
numéraire is a martingale}. Then, for a payoff $g$ paid at $T$:
\[
\frac{V_0}{N_0}=\E^N\!\left[\frac{g}{N_T}\right].
\]
The most convenient choice for options with a single expiry is $N_t=P(t,T)$, the zero-coupon bond maturing at $T$:
it is worth $P(0,T)=\DF$ today and $P(T,T)=1$ at expiry. The corresponding measure is called the \textbf{$T$-forward
measure} $\Q^T$, and the formula becomes
\begin{equation}\label{eq:prezzo-forward-misura}
\boxed{V_0=\DF\cdot\E^T[g]}
\end{equation}
Moreover, the forward price $F_t=F(t,T)$ is a \textbf{martingale} under $\Q^T$:
\begin{equation}
\E^T[F_T\mid\mathcal F_t]=F_t,\qquad F_T=S_T,\qquad \E^T[S_T]=F.
\end{equation}

\begin{perche}[title={Why is the forward a martingale under $\Q^T$?}]
Without dividends $F_t=S_t/P(t,T)$: it is the price of a security (the underlying) divided by the numéraire, hence a
martingale by the definition of $\Q^T$. With dividends the argument extends by using the portfolio that reinvests
the dividends. Intuitively: entering a forward costs nothing, so under the pricing measure its ``expected gain'' must
be zero.
\end{perche}

\begin{idea}[title={Key idea: working with the forward}]
Under $\Q^T$ the forward has no drift. Everything about rates, dividends and stock-borrow cost is packed into the two
numbers $F$ and $\DF$. This is why all the pricing ``kernels'' of the project take as input only $(F,\DF,T)$ and the
model parameters, and \emph{never} $r$ or $q$. The advantages:
\begin{itemize}
\item for SPX the dividends are discrete and uncertain, and the stock-borrow rate is not observable: assuming a wrong
yield $q$ would shift all the forward prices and distort the volatility surface;
\item $F$ and $\DF$ are obtained directly from the quotes via parity (Chapter~\ref{cap:dati}), so they are
consistent with the option prices;
\item the formulas become simpler (Chapter~\ref{cap:black}).
\end{itemize}
\end{idea}

\section{Implied volatility as the market's ``language''}
In Chapter~\ref{cap:black} we will see Black's formula: given $(F,K,T,\DF)$, the price depends on a single
unobservable parameter, the volatility $\sigma$, and is strictly increasing in it. So every price (that respects the
no-arbitrage bounds) corresponds to \emph{one and only one} $\sigma$: the \textbf{implied volatility}. Traders
exchange prices expressed in implied volatility because it is a number that can be compared across different strikes
and expiries. If the Black model were true, the implied volatility would be the same for every strike; in the data it
is not (\emph{smile} and \emph{skew}, Chapter~\ref{cap:heston}). Much of this document is about models able to
reproduce this fact.

\begin{riassunto}
\begin{itemize}
\item Payoffs: call $(S_T-K)^+$, put $(K-S_T)^+$. SPX options are European and cash-settled.
\item $\DF={}$today's price of 1 paid at $T$; $F$ = the delivery price that makes a forward worth zero; the forward
with price $K$ is worth $\DF\,(F-K)$.
\item Absence of arbitrage $\Rightarrow$ parity $C-P=\DF\,(F-K)$, bounds, monotonicity and convexity in $K$, with no
model.
\item Price $=\DF\cdot\E^T[\text{payoff}]$ under the $T$-forward measure, where $F_t$ is a martingale.
\item The project always reasons in terms of $(F,\DF,T)$ and never $r,q$: the pricing kernels never see a rate or a
dividend.
\end{itemize}
\end{riassunto}

\begin{domande}[unbreakable]
\begin{enumerate}
\item Prove put--call parity starting from the payoffs. \emph{Hint:} $(S-K)^+-(K-S)^+=S-K$.
\item Why does parity hold for SPX options and not for those on SPY? \emph{Hint:} SPY options are American; early
exercise breaks the equality of the payoffs, and only bounds remain.
\item In the binomial example, why does the real probability not matter? \emph{Hint:} the price is the cost of the
replication, which does not depend on how likely each scenario is.
\item What does it mean that $C$ must be convex in $K$? \emph{Hint:} \eqref{eq:breeden}: density $\ge0$.
\end{enumerate}
\end{domande}
